\Needspace{18\baselineskip}
\subsection{Varying the curve coefficients}
\label{subsec:weighted-curve-product-models}

The coefficient $f$ of the auxiliary curve in \cref{cor:nonclosed-separated-product-metric} does not affect the metric on the model locus. Scaling the curve forms themselves has a different effect: a positive coefficient rescales the polar distance in its coordinate, whereas a zero coefficient removes the coordinate altogether, so the diameter depends on the remaining curve areas. We describe the resulting family of model spaces and then compare its members.

Keep the product $M=Y\times C_0\times Z$, its projections and the forms of that corollary. For $a=(a_1,\ldots,a_m)\in[0,\infty)^m$ and a smooth function $f\geq0$ on $Y$, put
\[
 J(a)=\{i:a_i>0\},\qquad
 \beta_a=\sum_i a_i\pi_i^*\eta_i,\qquad
 \Theta_{a,f}=p^*\beta_a+(f\circ p)s^*\tau+r^*\kappa.
\]
The coordinates $i\in J(a)$ are called \emph{active}. For normalized ordinary models $q_i$ in the fixed curve backgrounds $\eta_i$, define
\[
 \Phi_a(q)=\Psi+\sum_{i\in J(a)}a_iq_i\circ\pi_i\circ p,
 \qquad \Psi=\psi\circ r.
\]
Terms with zero coefficient are omitted, also on their pole sets, so $\Phi_a(q)$ depends only on the active coordinates. Write $\mathcal S_a$ for this locus of potentials, and $q_{i,\mathcal I}=P_{\eta_i}[q_i]_{\mathcal I}$.

\Needspace{13\baselineskip}
\begin{lemma}\label{lem:weighted-product-model-projections}\index{Envelope!separated curve models}
For every $a$ and $f$ as above,
\[
 P_{\Theta_{a,f}}[\Phi_a(q)]=\Phi_a(q),\qquad
 P_{\Theta_{a,f}}[\Phi_a(q)]_{\mathcal I}=\Phi_a(q_{\mathcal I}).
\]
Both potentials have supremum zero. In particular, $\mathcal S_0=\{\Psi\}$.
\end{lemma}
\begin{proof}
For $a_i>0$, scaling the ordinary competitors shows that $a_iq_i$ is an ordinary model for $a_i\eta_i$. Its polar curvature is $a_i\mu_i$, where $\mu_i$ is the polar curvature of $q_i$. Taking atomic parts and applying \cref{thm:curve-background-homothety}, we see that its ideal model is $a_iq_{i,\mathcal I}$.

Put $w_a=\sum_{i\in J(a)}a_iq_i\circ\pi_i$. If there are at least two active factors, \cref{lem:separated-curve-model-projections} makes the corresponding sum an ordinary model on the product of the active curves, with its positive product background. For one active factor, this is the identity on the curve. The identity is preserved under pullback to $Y$, by \cref{prop:psh-descent-submersion}: every competitor in the pulled-back background descends, and so do its obstacle bound and its ordinary order, since the projection is surjective. If there are no active factors, the normalized model is zero by the maximum principle. Hence $P_{\beta_a}[w_a]=w_a$ in all cases.

On $B=Y\times C_0$, use the background $\beta_{a,f}=\pi^*\beta_a+(f\circ\pi)s_0^*\tau$. As in \cref{cor:nonclosed-separated-product-metric}, replace an ordinary competitor by its maximum with $w_a\circ\pi$ and restrict it to each section $Y\times\{z\}$. This gives $P_{\beta_{a,f}}[w_a\circ\pi]=w_a\circ\pi$.

For the ideal identity, take $g_a=\sum_{i\in J(a)}a_i g_i\circ\pi_i\circ\pi$, with the atomic Green potentials $g_i$ of \cref{lem:curve-atomic-ideal-model}. Scaling, addition and pullback preserve multiplier-ideal equivalence by \cref{lem:hermitian-i-calculus,lem:submersion-multiplier-pullback}. The complex Hessian of $g_a$ is a closed countable divisor current of finite mass minus a smooth form, and $g_a$ is $\beta_{a,f}$-plurisubharmonic. It is ordinarily equivalent to $w_{a,\mathcal I}\circ\pi$, where $w_{a,\mathcal I}=\sum_{i\in J(a)}a_iq_{i,\mathcal I}\circ\pi_i$. By \cref{lem:divisorial-potential-domination}, every ideal competitor is an ordinary competitor for $g_a$, so by the ordinary identity applied to the models $a_iq_{i,\mathcal I}$ it is bounded by $w_{a,\mathcal I}\circ\pi$. The latter is itself an ideal competitor, so the ideal envelope equals it. When $J(a)$ is empty, both envelopes on $B$ are zero directly. Finally, \cref{prop:fibration-obstacle-reduction}, applied to $M\to B$ and the divisor $B\times D$ with weight one, gives the identities upstairs, and independence of the factors gives normalization.
\end{proof}

For two tuples $q,\widetilde q$, let $\mu_i,\nu_i$ be their polar curvature measures in the fixed backgrounds $\eta_i$, and set
\begin{equation}\label{eq:weighted-product-polar-metric}
 \rho_a(\Phi_a(q),\Phi_a(\widetilde q))
       =\sum_{i\in J(a)}a_i\|\mu_i-\nu_i\|_{\mathrm{TV}}.
\end{equation}
On the singleton $\mathcal S_0$ this distance is zero.

\Needspace{20\baselineskip}
\begin{proposition}\label{prop:weighted-product-polar-metric}\index{Metric completeness!vanishing curve coefficients}\index{Model potential!weighted product collapse}
For every $a$, the locus $\mathcal S_a$ is a complete geodesic metric space for $\rho_a$, independently of $f$. Its ideal subspace is closed, geodesically convex and equal to the metric closure of its models with analytic singularities. Ideal projection is the unique nearest-point retraction onto this subspace and is $1$-Lipschitz, with distance
\[
 \rho_a\bigl(\Phi_a(q),\Phi_a(q_{\mathcal I})\bigr)
       =\sum_{i\in J(a)}a_i\mu_i^{\mathrm{na}}(C_i),
\]
where $\mu_i^{\mathrm{na}}$ is the nonatomic part of $\mu_i$. The diameter is $2\sum_i a_iA_i$, and metric convergence implies $L^1(M)$ convergence for every smooth positive volume form. The upper-volume comparison vanishes on this locus.

If $a,b$ have the same nonempty active set $J$, the correspondence $T_{ba}\Phi_a(q)=\Phi_b(q)$ is well defined, commutes with ideal projection and satisfies
\[
 \min_{i\in J}\frac{b_i}{a_i}\,\rho_a(u,v)
 \leq\rho_b(T_{ba}u,T_{ba}v)
 \leq\max_{i\in J}\frac{b_i}{a_i}\,\rho_a(u,v).
\]
For a fixed smooth positive volume form on $M$, there are constants $K_i$, independent of $a,b$ and of the normalized model tuple $q$, such that
\[
 \|\Phi_a(q)-\Phi_b(q)\|_{L^1(M)}
       \leq\sum_i K_i|a_i-b_i|.
\]
Consequently, as all $a_i$ tend to zero, the diameter tends to zero and every potential in $\mathcal S_a$ tends to $\Psi$ in $L^1(M)$, uniformly over the normalized model tuples.
\end{proposition}
\begin{proof}
Restriction to coordinate sections on which the other potentials are finite determines every active term $a_iq_i$ up to a constant, and normalization makes these constants zero. Thus \eqref{eq:weighted-product-polar-metric} is well defined, and it identifies $\mathcal S_a$ with the product of the spaces of curve models in the active coordinates, with their distances multiplied by $a_i$. Coordinates with $a_i=0$ do not appear, and for $J(a)=\varnothing$ the empty product is a single point.

The completeness, geodesic and atomic-projection results of \cref{cor:curve-model-completeness,prop:curve-analytic-distance-defect} apply coordinate by coordinate, and the weighted distances add; this proves the metric assertions and the distance formula. The diameter bound is attained by choosing polar measures $A_i\delta_{x_i}$ and $A_i\delta_{y_i}$ with $x_i\ne y_i$ on each active curve. Fiber integration and \cref{thm:curve-model-capacity-convergence} give $L^1$ convergence upstairs. The preceding lemma shows that ideal projection does not depend on $f$. Since every model and every pairwise maximum is $\Psi$ plus a pullback from $B$, \cref{prop:fibration-obstacle-reduction} gives zero upper volumes, so the upper-volume comparison vanishes.

For analytic density, approximate the atomic part $a_i\mu_i^{\mathrm{at}}$ of each scaled polar measure, where $\mu_i^{\mathrm{at}}=\mu_i-\mu_i^{\mathrm{na}}$, by finite atomic measures with rational coefficients and mass at most $a_iA_i$. These approximations converge in total variation. Their curve models in the background $a_i\eta_i$ have finitely many logarithmic poles with rational coefficients, whose denominators, together with the unit coefficient of $D$, can be cleared in local product monomials. They therefore give analytic models in $\mathcal S_a$ converging to every ideal model. Conversely, suppose that $\Phi_a(q)$ has analytic singularities. Restrict it to a section along an active curve on which all other potentials are finite and whose $Z$-coordinate lies outside $D$. By \cref{prop:analytic-type-calculus}, the restriction, which is $a_iq_i$ plus a constant and not identically $-\infty$, has analytic singularities. It is an ordinary model in the positive background $a_i\eta_i$, so \cref{prop:analytic-envelope-type} shows that it equals its ideal projection. The preceding lemma now shows that $\Phi_a(q)$ is an ideal model, and closedness gives the analytic-closure assertion. The argument uses that the scaled poles $a_iq_i$ themselves have rational coefficients; for irrational $a_i$, the scaled version of an analytic curve model need not be analytic.

For a common active set, the normalized tuple is determined by the potential. The two bounds for $T_{ba}$ follow by comparing each positive weight $b_i$ with $a_i$, and atomic projection in each coordinate commutes with this correspondence.

Finally, normalized compactness, \cref{thm:normalized-compactness}, bounds the $L^1(C_i)$ norms of all the $q_i$ uniformly. Integration along the compact fibers of $\pi_i\circ p$ gives constants $K_i$ bounding the $L^1(M)$ norms of their pullbacks. Off the pole sets, which have smooth volume zero, the difference of the two potentials is the sum of the functions $(a_i-b_i)q_i\circ\pi_i\circ p$. The triangle inequality proves the uniform estimate, also when some coefficients vanish. Taking $b=0$ gives uniform convergence to $\Psi$, and the diameter formula shows that the diameters tend to zero.
\end{proof}

Thus letting the auxiliary coefficient vanish leaves the metric of \cref{cor:nonclosed-separated-product-metric} unchanged, whereas letting a curve coefficient vanish removes its coordinate and decreases the diameter. These computations concern only the specified product loci, not general ordinary models in higher dimension.

\Needspace{18\baselineskip}
\subsubsection{Comparing bounded model spaces}
\label{subsec:weighted-product-metric-comparison}

When the active sets differ, there is no natural map between the spaces $\mathcal S_a$, and we compare them as abstract metric spaces. They are bounded, but they are not compact as soon as some coefficient is positive. The Gromov--Hausdorff distance $d_{\mathrm{GH}}$ of bounded metric spaces from \cref{def:bounded-gh-comparison}, together with its estimate through correspondences in \cref{lem:bounded-metric-correspondence}, applies without a compactness hypothesis.

Write $\mathcal I_a\subset\mathcal S_a$ for the ideal subspace of \cref{prop:weighted-product-polar-metric}, with the restricted metric $\rho_a$.

\Needspace{19\baselineskip}
\begin{proposition}\label{prop:weighted-product-metric-comparison}\index{Model potential!coefficient comparison}
For all nonnegative coefficient vectors $a,b$, including vectors with different active sets,
\begin{equation}\label{eq:weighted-product-gh-bounds}
 \left|\sum_i(a_i-b_i)A_i\right|
 \leq d_{\mathrm{GH}}(\mathcal S_a,\mathcal S_b)
 \leq\sum_i|a_i-b_i|A_i.
\end{equation}
The same bounds hold for $d_{\mathrm{GH}}(\mathcal I_a,\mathcal I_b)$. If $b_i\leq a_i$ for every $i$, both distances equal $\sum_i(a_i-b_i)A_i$. In particular, the distance from $\mathcal S_a$ to the singleton $\mathcal S_0=\mathcal I_0=\{\Psi\}$ equals $\sum_i a_iA_i$.

For $b\leq a$, the map $T_{ba}\Phi_a(q)=\Phi_b(q)$ is a surjective $1$-Lipschitz map and commutes with ideal projection. Every fiber has diameter
\[
 2\sum_{\{i:a_i>0,\ b_i=0\}}a_iA_i.
\]
If some $a_i>0$, then $\mathcal S_a$ and $\mathcal I_a$ are not totally bounded, and hence not compact.
\end{proposition}
\begin{proof}
Use the correspondence
\[
 R_{ab}=\{(\Phi_a(q),\Phi_b(q)):q_i
                \text{ are normalized ordinary curve models}\}.
\]
Every potential in either locus arises from such a tuple, with arbitrary normalized models in its coordinates of zero coefficient, so both projections of $R_{ab}$ are surjective. For any two tuples, put $d_i=\|\mu_i-\nu_i\|_{\mathrm{TV}}$. By \eqref{eq:weighted-product-polar-metric}, the difference of their two distances is $\sum_i(a_i-b_i)d_i$, with $0\leq d_i\leq2A_i$. Thus $\operatorname{dis}R_{ab}\leq2\sum_i|a_i-b_i|A_i$, and \cref{lem:bounded-metric-correspondence} gives the upper bound in \eqref{eq:weighted-product-gh-bounds}. The diameter formula in \cref{prop:weighted-product-polar-metric} and the same lemma give the lower bound.

Restricting the correspondence to tuples of ideal curve models, we obtain a correspondence between the two ideal subspaces with the same distortion bound. The atomic measures $A_i\delta_{x_i}$ and $A_i\delta_{y_i}$ show that the ideal subspaces have the same diameters as the full spaces, so the lower bound also holds for them. When $b\leq a$, the upper and lower bounds coincide, including when $b=0$.

For the transition map, every active coordinate of $b$ is active for $a$, and uniqueness of the normalized active tuple makes the map well defined. Every $b$-tuple extends to an $a$-tuple, so the map is surjective, and $\sum_i b_i d_i\leq\sum_i a_i d_i$ proves the Lipschitz bound. Atomic projection in each coordinate commutes with discarding and rescaling coordinates. In a fiber, every coordinate with $b_i>0$ is fixed, while the coordinates with $a_i>0$ and $b_i=0$ vary freely. The diameter of the fiber is therefore the sum above, by the calculation of the diameter on each curve.

Finally, fix an active coordinate $i$ and all other coordinates at their zero models. Let $q_i$ have polar measure $A_i\delta_x$, with $x$ ranging over infinitely many distinct points of $C_i$. The resulting potentials are ideal models, and any two of them are at distance $2a_iA_i$. This infinite separated set excludes total boundedness of both spaces. In particular, the comparison above is between complete bounded spaces that are not compact.
\end{proof}
